%=====================================================================
\section{Discussion}\label{sec:discussion}
%=====================================================================

\subsection{Principal findings}\label{sec:disc_findings}
\paragraph{RQ1: learning from heterogeneous annotation.}
Treating unannotated entries as unknown rather than negative is the single most effective change for seen diagnoses: masked loss improves on naive pooling by 0.7 points of macro AUROC and 0.6 points of AUPRC (Table~\ref{tab:loso}), and in the controlled experiment the masked objective keeps retained diagnoses within 2.8 points of their complete-annotation accuracy while half of the annotation is removed (Table~\ref{tab:e2}). This is the behaviour predicted by Proposition~\ref{prop:naive}. On top of correct label handling, the hybrid head adds a further, small but consistent improvement: it has the highest mean AUROC and AUPRC of all methods, is better than every baseline in at least eight of twelve fold--seed pairs, and the difference is significant against six of the seven baselines. It is not significant against the masked baseline, and we do not claim that the hybrid head improves seen-diagnosis accuracy over well-implemented masking; the defensible statement is that it matches the best label-handling baseline on seen diagnoses while adding capabilities that no baseline has.

\paragraph{RQ2: transfer to an unseen database.}
None of the four domain-generalization methods improved on masking alone, and DANN was the weakest method after naive pooling. This agrees with the observation that DG penalties tuned for one benchmark rarely help consistently across held-out domains~\cite{ballas2022}, and it suggests that, in pooled public ECG data, the label-induced shortcut of Proposition~\ref{prop:naive} is a larger source of transfer error than the feature-statistics differences that DG methods target. Within each held-out database, the hybrid model is significantly better than most baselines on PTB-XL, Georgia and Shaoxing--Ningbo, but not on CPSC, whose small and rhythm-centred vocabulary leaves the least room for a text head that shares information between related diagnoses. Weight averaging, which we applied to all methods, improved every baseline by 0.3--1.8 points (Appendix~\ref{secA2}); it is an inexpensive component that published multi-source ECG comparisons should include in all arms.

\paragraph{RQ3: scoring diagnoses without labels.}
The clearest result of the study concerns diagnoses for which no training label exists. On the five Shaoxing--Ningbo diagnoses that no training database annotates, the text head reaches a macro AUROC of 67.0, against chance for every baseline, and on bradycardia, which the training data supervise only through an implication rule, it reaches 94.3 against 89.8 for the best linear head. In the controlled experiment, removed PTB-XL diagnoses are scored at 66.6--69.0 by the text head and at or below chance by the baselines, and per-diagnosis results show near-complete recovery of conduction and rhythm findings whose descriptions are close to those of supervised diagnoses (right bundle branch block 93.5, left anterior fascicular block 94.4, atrial flutter 87.1). The systematic evaluation E3 shows that this is not a property of a few favourable diagnoses: over 120 withheld diagnosis--database pairs, the text head scores 88\% above chance with a macro AUROC of 70.2, retaining on average 63\% of the above-chance AUROC that full supervision reaches. These are not clinically adequate accuracies for most diagnoses, but they are obtained without a single label, from a model whose seen-diagnosis accuracy is unchanged. Zero-shot accuracy increases with the semantic proximity between the description and the nearest trained diagnosis ($\rho=0.26$), and the failures have a consistent cause: the worst-scored diagnoses (sinus tachycardia, sinus bradycardia, ST elevation and depression, right axis deviation, low QRS voltages) have a close neighbour whose description shares the topic but differs in a number or a sign (``a rate above 100 beats per minute'' against ``below 60''; ``elevation'' against ``depression''), and they are scored below chance. A frozen sentence encoder represents what a criterion is about but not reliably what it specifies, which identifies descriptions with explicit, encoder-visible contrasts, or encoders trained on ECG criteria, as the most direct route to better zero-shot accuracy.

\subsection{Why a hybrid head}\label{sec:disc_hybrid}
Our initial design, a text-only label-query head trained with teacher pseudo-labels (Appendix~\ref{secA2}), scored unseen diagnoses well but was 1.1 points of mean AUROC below the Li-style baseline on seen diagnoses in matched single-seed runs, and its validation AUROC dropped abruptly, by 7 to 31 points, on all four folds once pseudo-labels took effect (Fig.~\ref{fig:curves}). The diagnosis of that failure motivated the final design. A text-only head must represent every diagnosis through a shared projection of its description, which is an advantage for diagnoses without labels and a constraint for diagnoses with thousands of labels, where a free weight vector is more accurate. The hybrid resolves the trade-off without a hyper-parameter: seen diagnoses keep a free weight vector and gain an ensemble partner, unseen diagnoses use the only head that can score them, and the partition between the two cases is known from the training annotation. The ablations support this reading. The gain on seen diagnoses comes from fusion rather than from auxiliary training, joint training of the encoder adds to a post-hoc ensemble, and the text head is better at unseen diagnoses when trained alongside a linear head (67.0) than alone (62.0), presumably because the linear head shapes the shared encoder towards features that are discriminative per diagnosis. The cost is 1.27\,M parameters and two to three minutes of additional training per run.

\subsection{Relation to previous work}\label{sec:disc_related}
Partial-label methods for ECG~\cite{li2024} and for images~\cite{durand2019,benbaruch2022} correspond to our masked and Li-style baselines, which we found to be strong; our results extend their evaluation from in-distribution test splits to held-out databases and show that their benefit persists there. ECG--language models pretrained on report-paired data~\cite{merl2024,suprime2025,melp2025} obtain open-vocabulary scoring from hundreds of thousands of reports; our approach obtains a weaker but non-trivial form of it from the categorical labels that public databases already carry, by modeling which labels each source observes, on a single cloud GPU and in about twelve minutes per run. Our direct comparison with MERL quantifies the trade-off. On 120 withheld common diagnoses, report pretraining is significantly better on average (77.3 against 70.2), mainly on rhythm diagnoses that reports describe very often; the text-anchored head is better on 45 of the 120 pairs, particularly on repolarization and morphology findings, and on the diagnoses that are genuinely rare and absent from all our training data (67.0 against 52.9--62.6 on Shaoxing--Ningbo; 94.3 against 86.0 on bradycardia). On the seen diagnoses, supervised training on the public labels (84--88) is far ahead of zero-shot report pretraining (76--78). The two sources of knowledge are therefore complementary rather than ordered: a report-pretrained encoder could replace our randomly initialized encoder, and our annotation-aware objective and hybrid head could be used to fine-tune such a model on heterogeneously annotated data, which we consider the most promising extension of this work.

\subsection{Clinical implications}\label{sec:disc_clinical}
Two aspects of the design bear directly on deployment, independently of the magnitude of the measured effects. First, Proposition~\ref{prop:naive} shows that a model trained by naive pooling estimates $P(y=1\mid\mathbf{x})$ multiplied by the probability that the recording comes from a source that annotates the diagnosis. In a hospital that did not contribute training data this factor is extrapolated from features that identify the training sources and has no meaning for the new site, so the scores are deflated by an amount that varies across diagnoses; the masked objective removes this factor by construction. It does not, however, make the probabilities calibrated: all methods, including ours, over-predict by a factor of about three because of the shared asymmetric loss (Section~\ref{sec:res_calibration}), so probabilities must be recalibrated locally before they are read as risks. Second, because the text head contains no diagnosis-specific weights, a site that wishes to screen for a finding that it reports but that no public database codes can do so by writing a description, without collecting labels or retraining. The accuracies we observe for such diagnoses (macro AUROC 70 over 120 withheld diagnoses, above 85 for close relatives of trained diagnoses such as bundle branch and atrioventricular blocks, and below chance for rate- and polarity-defined findings) make this a tool for triage, retrospective cohort search or prioritizing records for expert annotation, not for diagnosis; any new description should be validated on local labeled data before use. Finally, the gap of about eight points between macro F1 at source-tuned and at oracle thresholds, similar for every method, shows that decision thresholds transfer poorly across databases. Deployment at a new site therefore requires local threshold calibration on a small labeled sample regardless of the training method, and threshold-free metrics such as AUROC overstate what a fixed-threshold system would achieve out of the box.

\subsection{Limitations}\label{sec:limitations}
Several limitations qualify the results. The improvement on seen diagnoses is small (0.25--0.97 points of mean AUROC) and not significant against the masked baseline; it was also not anticipated by the selection data, on which the selected configuration was 0.0004 below the strongest baseline (Section~\ref{sec:selection}). It is consistent across folds and seeds but is an improvement of the order that architectural refinements report on in-distribution benchmarks, and its clinical relevance on seen diagnoses is limited. Genuinely unseen diagnoses are few (six in two folds); E3 compensates by withholding 120 diagnosis--database pairs, but a withheld common diagnosis is not the same as a rare finding that no database codes, and its related diagnoses remain in training. The annotation-removal experiment uses a single seed and a single database. The comparison with MERL uses one released checkpoint; its own prompts were mapped by hand to our vocabulary, its checkpoint was selected on splits of three of our held-out databases, and other ECG--language models~\cite{suprime2025,melp2025} were not evaluated. The annotation threshold was varied only between 5 and 20 positives and with one seed. Signals are resampled to 100\,Hz, which preserves the diagnoses in our vocabulary but removes high-frequency content relevant to findings such as pacing spikes or fragmented QRS, and places all methods at the same disadvantage. Patient identifiers are not included in the Challenge redistribution, so the E1 folds are separated by database, which is the intended unit of analysis, while the E2 split on PTB-XL is by record and may be slightly optimistic for all methods. Labels refer to whole recordings, whereas training uses a random 10\,s window and inference the first 20\,s; for long CPSC recordings a window may miss the paroxysmal event that justified the label. Shaoxing--Ningbo pools two hospitals and CPSC two releases of one challenge, so a fold does not always correspond to a single institution. The diagnosis descriptions were written by the authors from published recommendations rather than by an independent panel, and the open-vocabulary results depend on their wording, as the name-only ablation in Appendix~\ref{secA2} shows. The DG loss weights were fixed a priori and not tuned, so the DG baselines may be under-tuned; the hyper-parameters of the proposed method were also fixed a priori, except for the choice among four hybrid variants on one validation fold. All data are retrospective public releases whose labeling practices we cannot audit; probabilities are miscalibrated for all methods and were not recalibrated, and subgroup performance by age and sex was not assessed.

%=====================================================================
\section{Conclusion}\label{sec:conclusion}
%=====================================================================
We studied 12-lead ECG classification from several public databases whose annotation vocabularies differ, formalized why pooling them naively trains a model to recognize the source, and proposed a hybrid classifier in which a conventional linear head and a text-anchored label-query head share one encoder, are trained with a loss restricted to the labels each source observes, and are fused on diagnoses that the training data annotate. Under a leave-one-database-out protocol over four databases from three countries, with three seeds and multiplicity-corrected tests, the method matched the strongest label-handling baseline on seen diagnoses and was significantly better than naive pooling and all four domain-generalization baselines, and it was the only one of the trained methods able to score diagnoses for which no training label existed: 67.0 macro AUROC on five such diagnoses where all baselines were at chance, 66.6--69.0 on diagnoses removed from training in a controlled experiment, and 70.2 over 120 systematically withheld diagnosis--database pairs, 88\% of them above chance. A report-pretrained ECG--language model was better on average on the withheld common diagnoses (77.3) but worse on the genuinely rare ones, and the failures of the text head were traced to descriptions that differ from a trained neighbour only in a number or a sign. Correct handling of missing annotation mattered more for cross-database transfer than any of the domain-generalization penalties we evaluated. Future work includes initializing the encoder from report-pretrained ECG--language models, descriptions written and validated by independent cardiologists with explicit numerical criteria, patient-level evaluation where identifiers exist, calibration analysis, and prospective evaluation at a site that contributed no training data.
